\documentstyle[11pt]{article}
\setlength{\textwidth}{170mm}
\setlength{\textheight}{220mm}
\setlength{\oddsidemargin}{0mm}
\setlength{\evensidemargin}{10mm}
\setlength{\topmargin}{-10mm}
\def\RR{\hbox{\raise.4ex\hbox{$\scriptstyle{|}$}\hskip-1.85pt R}}
\newcommand{\REF}{\nopagebreak\noindent{\large{\bf
      References}}\vspace{.2in}}

\newcommand{\Ref}[1]{\medskip\noindent\hangafter=1
      \hangindent=20pt\ignorespaces #1}
\pagestyle{empty}
\begin{document}


\noindent
University of Illinois \hfill Department of Economics \\
Spring 1997 \hfill Econ 476 \hfill Roger Koenker

\vspace{4mm}
\begin{center}

\large{
{\bf Some Bibliography }}
\end{center}

\vspace{10mm}
Inevitably, one of the first questions raised by students in Econ 476 is, ``Can you suggest some other references which might be helpful for 
background reading?''  I am usually reluctant to do this because there is, 
inevitably, some question about guessing an appropriate level for the 
supplementary material.  To circumvent this problem I have tried to 
assemble some references which span a reasonably wide range of levels.  The
references are subdivided into the two traditional categories:  Probability 
and Statistics.  Within these categories they are listed in order of 
approximate difficulty.

A useful overview of some of this material is contained in four books
written explicitly by and for econometricians again listed in order of
difficulty.

\vspace{2mm}
\Ref 1.   McCabe, B. and Tremayne, A. (1993), {\it Elements of Modern
Asymptotic Theory with Statistical Applications}, Manchester U. Press.

\Ref 2.  White, H. (1984) {\it Asymptotic Theory for Econometricians},
Academic Press.

\Ref 3. White, H. (1994) 
{\it Estimation, Inference, and Specification Analysis}, 
Cambridge University Press.

\Ref 4.  Davidson, J. (1994) {\it Stochastic Limit Theory}, Oxford U. Press.

\vspace{5mm}
\noindent
Probability

\vspace{3mm}
\Ref 1.   Chung, K.L. (1979), {\it Elementary Probability Theory}, Springer-Verlag.

\Ref 2.   Whittle, P. (1970), {\it Probability via Expectation}, Springer-Verlag.

\Ref 3.   Breiman, L. (1968), {\it Probability} (Siam reprint 1992).

\Ref 4.   Chung, K.L. (1974), {\it A Course in Probability Theory}, Academic Press.

\Ref 5.   Williams, D. (1991), {\it Probability with Martingales}, Cambridge.

\Ref 6.   Billingsley, P. (1979), {\it Probability and Measure}, Wiley.

\Ref 7.   Pollard, D. (1984), {\it Convergence of Stochastic Processes}, Springer-Verlag.

\vspace{3mm}
\noindent
Statistics

\vspace{3mm}
\Ref 1.  Bickel, P. and K. Doksum (1977), {\it Mathematical Statistics}, Holden-Day.

\Ref 2.  Welsh, A. H. (1996), {\it Aspects of Statistical Inference}, Wiley.

\Ref 3.   Cox, D.R. and D.V. Hinkley (1974), {\it Theoretical Statistics}, Chapman-Hall.

\Ref 4.   McCullagh,  and P. Nelder, J.A. (1989), {\it Generalized Linear Models}, Chapman-Hall.

\Ref 5.   Serfling, R. J. (1980), {\it Approximation Theorems of Mathematical 
Statistics}, Wiley.

\Ref 6.   Lehmann, E. (1983), {\it Theory of Point Estimation}, Wiley.

\Ref 7.   Hajek, J. and Sidak, Z. (1967), {\it Theory of Rank Tests}, Academic Press.

\end{document}
